\section{The remaining Yang--Mills obligations}\label{sec:open}

The results above do not establish the hypotheses of a continuum
mass-gap theorem. The following are separate missing implications.

\begin{enumerate}
\item \textbf{Collective nonlinear coercivity.}
The finite-block bound must be replaced or supplemented by an estimate
for an entire elimination layer in the actual compact-link vacuum,
uniform in volume and ultraviolet cutoff. Its validity for each fixed
block does not imply such a collective bound.

\item \textbf{Control after exact elimination.}
The estimates must survive the generated interactions, the new vacuum,
the induced coarse metric and the appropriate form domains.
The reduced operator has not been identified with another Hamiltonian
of the bare Wilson form \eqref{eq:hamiltonian}. Substitution of a running
coupling into the strong-coupling corollary is therefore insufficient.

\item \textbf{A physical terminal gap.}
At a fixed physical coarse scale, one must control the actual terminal
operator and its metric. The explicit bound at large bare $g^2$ provides
a possible comparison target, but no comparison from the actual
renormalized operator to that target is proved here.

\item \textbf{A nontrivial continuum field theory and spectral transfer.}
One must construct compatible renormalized observables, their limiting
correlations and a field theory with the requisite positivity, locality,
covariance and Yang--Mills short-distance behavior. A positive finite
physical gap requires an estimate uniform along the relevant limit.
There must also be a nonzero finite-energy sector: exclusion of spectrum
from an interval in a vacuum-only Hilbert space would not solve the
problem posed in~\cite{JW}.
\end{enumerate}

The scale distinction is especially important. The strong-coupling
bound is in lattice units. At a fixed admissible $x$, its physical lower
bound is proportional to $x/a$ and diverges as $a\downarrow0$.
Along a weak bare-coupling continuum trajectory, $x(a)$ instead tends
to zero, outside the theorem's certified window. Neither observation
disproves a continuum mass gap; both prevent an inference of such a
gap from the fixed-coupling estimate.

There is also a concrete obstruction to treating every configuration
as a positive quadratic fluctuation problem. On an even periodic
cubic lattice in spatial dimension at least two, choose central
$\SU(2)$ links
\[
 U_i(z)=(-I)^{\sum_{k<i}z_k}.
\]
Every elementary plaquette is $-I$. For a real link field $v$, vary
$U_i(z;t)=U_i(z)\exp(itv_i(z)\sigma_3/2)$. Then
\begin{equation}\label{eq:negative-hessian}
 \left.\frac{\mathrm d^2}{\mathrm dt^2}
 \frac{b}{x}\sum_p\left(1-\frac{\operatorname{Re}\Tr U_p(t)}2\right)
 \right|_{t=0}
 =-\frac{b}{4x}\sum_p|(\mathrm dv)_p|^2 .
\end{equation}
A circulation around one square has zero divergence, zero block mean
in each link component and nonzero curl. At this central background the
adjoint transports are the identity, so these are also the covariant
conditions. It is a discarded direction for the corresponding
componentwise block average and has strictly negative magnetic
Hessian for $b>0$. Thus an all-background positive-Hessian comparison
cannot provide the required nonlinear estimate. A negative Hessian of
the potential does not imply a negative excitation Hamiltonian or a
small quantum gap. Large-field regions require their own estimates.

The infinite-volume issue is kept distinct from uniformity in finite
volume. The general stability theorem proved here supplies a common
finite-volume lower bound. We do not claim in this paper to construct
the infinite-volume state and its GNS Hamiltonian at the explicit
threshold $59.9$. The earlier application of Yarotsky's results
addresses that construction in a sufficiently strong, nonexplicit
window~\cite{Yarotsky,Spectral}.

\section{Evidence, reproducibility and attribution}\label{sec:evidence}

The supplementary source archive contains the full \LaTeX{} project,
a standard-library Python verifier, exact arithmetic results, selected
historical proof notes, and selected Lean source certificates with their
toolchain manifest. The principal analytic proofs are in this paper.
The machine checks have the following narrower scope:
\begin{itemize}
\item rational upper bounds for the square-root sector constants,
the majorant supersolution inequalities and safe coupling thresholds;
\item the finite counting correction and its three-qubit matrix element;
\item exact finite-dimensional Schur and oscillator arithmetic.
\end{itemize}
The Python verifier uses rational arithmetic for its acceptance
conditions; decimal displays are explanatory only. No random search,
floating-point eigensolver, or finite sample is used to infer an
all-volume theorem. The optional Lean certificates concern scalar
inequalities, finite sets and a finite matrix calculation. They do not
formalize the analytic stability theorem, the Feynman--Kac construction,
Li--Yau theory or continuum quantum field theory.

The archive's \path{README.md} gives build and verification commands,
and \path{MANIFEST.json} records the included files and their SHA-256
hashes. Historical notes are included for provenance and may retain
their original notation or approximate threshold displays. The
statements and safe sufficient thresholds in the present manuscript
govern this release. In particular, $73.4$, rather than a downward
rounded $73.3$, is a safe decimal upper rounding of the displayed
four-spatial-dimensional link threshold when that comparison is used.
No claim of global novelty is inferred from the absence of an identified
prior formula.

\subsection*{AI assistance and review status}

OpenAI Codex assisted with source inspection, drafting, mathematical
consistency checks, exact verification code, typesetting and release
preparation. Separate AI-assisted readings were used to check the
mathematical arguments and the scope of the claims. This is not
independent human peer review, and it does not imply that every
derivation has been manually checked by the named author. The release
includes a review record stating the performed checks and their
limitations. No AI system is listed as an author.

\subsection*{License}

This manuscript and its \LaTeX{} source are licensed under Creative
Commons Attribution 4.0 International
(\href{https://creativecommons.org/licenses/by/4.0/}{CC BY 4.0}).
The accompanying repository-derived software and formal proof sources
retain their GNU Affero General Public License, version 3 or later.
The manuscript license does not relicense that software.
